% GENERATED FILE. Edit canonical lessons, then run npm run generate:books.
% canonical-source-hash: fnv1a64:6631128cffcf0cec
% canonical-lessons: TE-C136-jamapandu, TE-C136-pancadara, TE-C136-podavu, TE-C136-potti, TE-C136-baruvu

\chapter{Guava, Sugar, Tall, Short, Heavy}
\label{ch:te-guava-sugar-tall-short-heavy}
\hlchaptermodality{136}

\begin{chapteropening}
\textbf{By the end of this chapter:} \emph{I can say a guava, sugar, tall, short and heavy.}

Complete the last lesson of chapter 136: i can say a guava, sugar, tall, short and heavy.
\end{chapteropening}

\section[\texorpdfstring{\te{జామపండు} (jāmapaṇḍu)}{జామపండు (jāmapaṇḍu)}]{\te{జామపండు} (jāmapaṇḍu) — a guava}

\label{lesson:TE-C136-jamapandu}



\begin{quote}
Before the new one: say the Telugu for a lemon, then the Telugu for grapes.
\end{quote}

\subsection*{You'll want to know: \te{జామపండు}}
\textbf{\te{జామపండు}} — \emph{jāmapaṇḍu} — \textquotedblleft{}a guava\textquotedblright{}.

\begin{grammarlens}[title={What you've built}]
One more word for this chapter.
\end{grammarlens}

\subsection*{Your turn}
\begin{itemize}
\raggedright
  \item \emph{Say it:} \emph{jāmapaṇḍu}
  \item \emph{Say it:} \emph{jāmapaṇḍu}, once more
  \item \emph{Say it:} say \emph{jāmapaṇḍu}
  \item \emph{Recall:} say the Telugu for a lemon, then the Telugu for grapes, then say \emph{jāmapaṇḍu} again
\end{itemize}

\begin{tcolorbox}[breakable,colback=teal!4,colframe=teal!35!black,title={Before you move on}]
What does \te{జామపండు} mean? (\textquotedblleft{}A guava\textquotedblright{}.) Say it once more.
\end{tcolorbox}
\section[\texorpdfstring{\te{పంచదార} (pañcadāra)}{పంచదార (pañcadāra)}]{\te{పంచదార} (pañcadāra) — sugar}

\label{lesson:TE-C136-pancadara}



\begin{quote}
Before the new one: say the Telugu for grapes, then the Telugu for a guava.
\end{quote}

\subsection*{You'll want to know: \te{పంచదార}}
\textbf{\te{పంచదార}} — \emph{pañcadāra} — \textquotedblleft{}sugar\textquotedblright{}.

\begin{grammarlens}[title={What you've built}]
One more word for this chapter.
\end{grammarlens}

\subsection*{Your turn}
\begin{itemize}
\raggedright
  \item \emph{Say it:} \emph{pañcadāra}
  \item \emph{Say it:} \emph{pañcadāra}, once more
  \item \emph{Say it:} say \emph{pañcadāra}
  \item \emph{Recall:} say the Telugu for grapes, then the Telugu for a guava, then say \emph{pañcadāra} again
\end{itemize}

\begin{tcolorbox}[breakable,colback=teal!4,colframe=teal!35!black,title={Before you move on}]
What does \te{పంచదార} mean? (\textquotedblleft{}Sugar\textquotedblright{}.) Say it once more.
\end{tcolorbox}
\section[\texorpdfstring{\te{పొడవు} (poḍavu)}{పొడవు (poḍavu)}]{\te{పొడవు} (poḍavu) — tall, long}

\label{lesson:TE-C136-podavu}



\begin{quote}
Before the new one: say the Telugu for a guava, then the Telugu for sugar.

An adjective goes before the noun and does not change: \textbf{\te{పాత} \te{పుస్తకం}}, \textbf{\te{కొత్త} \te{పుస్తకం}}.
\end{quote}

\subsection*{You'll want to know: \te{పొడవు}}
\textbf{\te{పొడవు}} — \emph{poḍavu} — \textquotedblleft{}tall, long\textquotedblright{}.

\begin{grammarlens}[title={What you've built}]
One more word for this chapter.
\end{grammarlens}

\subsection*{Your turn}
\begin{itemize}
\raggedright
  \item \emph{Say it:} \emph{poḍavu}
  \item \emph{Say it:} \emph{poḍavu}, once more
  \item \emph{Say it:} say \emph{poḍavu}
  \item \emph{Recall:} say the Telugu for a guava, then the Telugu for sugar, then say \emph{poḍavu} again
\end{itemize}

\begin{tcolorbox}[breakable,colback=teal!4,colframe=teal!35!black,title={Before you move on}]
What does \te{పొడవు} mean? (\textquotedblleft{}Tall, long\textquotedblright{}.) Say it once more.
\end{tcolorbox}
\section[\texorpdfstring{\te{పొట్టి} (poṭṭi)}{పొట్టి (poṭṭi)}]{\te{పొట్టి} (poṭṭi) — short}

\label{lesson:TE-C136-potti}



\begin{quote}
Before the new one: say the Telugu for sugar, then the Telugu for tall.
\end{quote}

\subsection*{You'll want to know: \te{పొట్టి}}
\textbf{\te{పొట్టి}} — \emph{poṭṭi} — \textquotedblleft{}short\textquotedblright{}.

\begin{grammarlens}[title={What you've built}]
One more word for this chapter.
\end{grammarlens}

\subsection*{Your turn}
\begin{itemize}
\raggedright
  \item \emph{Say it:} \emph{poṭṭi}
  \item \emph{Say it:} \emph{poṭṭi}, once more
  \item \emph{Say it:} say \emph{poṭṭi}
  \item \emph{Recall:} say the Telugu for sugar, then the Telugu for tall, then say \emph{poṭṭi} again
\end{itemize}

\begin{tcolorbox}[breakable,colback=teal!4,colframe=teal!35!black,title={Before you move on}]
What does \te{పొట్టి} mean? (\textquotedblleft{}Short\textquotedblright{}.) Say it once more.
\end{tcolorbox}
\section[\texorpdfstring{\te{బరువు} (baruvu)}{బరువు (baruvu)}]{\te{బరువు} (baruvu) — heavy}

\label{lesson:TE-C136-baruvu}



\begin{quote}
Before the new one: say the Telugu for tall, then the Telugu for short.
\end{quote}

\subsection*{You'll want to know: \te{బరువు}}
\textbf{\te{బరువు}} — \emph{baruvu} — \textquotedblleft{}heavy\textquotedblright{}.

\begin{grammarlens}[title={What you've built}]
One more word for this chapter.
\end{grammarlens}

\subsection*{Your turn}
\begin{itemize}
\raggedright
  \item \emph{Say it:} \emph{baruvu}
  \item \emph{Say it:} \emph{baruvu}, once more
  \item \emph{Say it:} say \emph{baruvu}
  \item \emph{Recall:} say the Telugu for tall, then the Telugu for short, then say \emph{baruvu} again
\end{itemize}

\begin{tcolorbox}[breakable,colback=teal!4,colframe=teal!35!black,title={Before you move on}]
What does \te{బరువు} mean? (\textquotedblleft{}Heavy\textquotedblright{}.) Say it once more.
\end{tcolorbox}
